\section*{Scope and conventions}
\addcontentsline{toc}{section}{Scope and conventions}
We work with small categories relative to a fixed Grothendieck universe, allowing a larger universe for their category. We write $BA=|NA|$. Every space is represented by a CW complex or a simplicial set with its derived homotopy type. The homotopy cofiber of an unpointed map $f:X\to Y$ is the pointed homotopy pushout $C(f)=Y\amalg_X^h *$, equivalently the cofiber of $f_+:X_+\to Y_+$ in pointed spaces. In particular $C(\varnothing\to Y)=Y_+$.

The word fundamental follows the prompt's relative-Theorem-A convention. A weak localizer in Chapter XVI means only Loc (1)--(3); the two meanings must not be interchanged. The symbols $R$, $T$, $W_D$, $W_D^{\mathrm{lc}}$, and $W_D^1$ retain the prompt's definitions. An $RT$-fixed localizer is called a homology localizer here; we do not identify all localizing subcategories of spectra with smash-homology kernels.

The attached LaTeX files are transcriptions, not facsimiles. We can check their internal mathematics but cannot certify uncertain manuscript glyphs without the scans. An inference about a reading is explicitly marked as a mathematical inference. Check events and source access are recorded separately from proofs.
